\documentclass[11pt]{article}
\usepackage{fontspec}
\usepackage[french,english]{babel}
\usepackage{amsmath}
\usepackage{amssymb}
\usepackage{amsthm}
\usepackage{mathrsfs}
\usepackage{tikz-cd}
\usepackage{enumitem}

\newcommand{\printedpage}[1]{\marginpar{\footnotesize p.~#1}}
\newcommand{\msp}[1]{\marginpar{\footnotesize [#1]}}
\newcommand{\Hom}{\operatorname{Hom}}
\newcommand{\uHom}{\underline{\operatorname{Hom}}}
\newcommand{\Ob}{\operatorname{Ob}}
\newcommand{\id}{\mathrm{id}}
\newcommand{\bDelta}{\boldsymbol{\Delta}}

\setlength{\parindent}{0pt}

\begin{document}

\printedpage{4}
% head: A. Grothendieck : Dérivateurs, Ch. IV, Édition des Universités de Montpellier II et Paris VII
% foot: Version 2 février 2002
\msp{page 78}
a ``canonical'', ``natural'', ``willed by God'' way of ``fleshing it out'' into a substantial diagram. The most obvious arrows that present themselves are relative to a functor
\[ f : I \longrightarrow J \]
in Diag, and to the two adjunction morphisms for the pair $(f^*, f_*)$, namely
\[
\begin{aligned}
G &\longrightarrow f_*f^*G && \text{in } D(J)\\
f^*f_*F &\longrightarrow F && \text{in } D(I). \quad (\text{\footnotemark})
\end{aligned}
\]
\footnotetext{Of course, dual problem for $(f_!, f^*)$.}

I am going to try to flesh out these two arrows, and am going to generalize by considering the situation opposite \texttt{[below]}
\[
\begin{tikzcd}
I \arrow[d, "f"'] & F \\
J & G, \arrow[u, "u"']
\end{tikzcd}
\]
where
\[ \Hom^f(G,F) \overset{\text{def}}{=} \Hom(f^*G,F) \overset{\text{adj.}}{\simeq} \Hom(G,f_*F). \]
Thus the same $u$ is interpreted by \textit{two} skeletal arrows
\[
\left\{
\begin{aligned}
f^*(G) &\longrightarrow F && \text{in } \mathbf{D}(I)\\
G &\longrightarrow f_*(F) && \text{in } \mathbf{D}(J).
\end{aligned}
\right.
\tag*{(1)}
\]
To flesh them out is to interpret them (if it can be done) in terms of two objects (canonical if possible) in $\mathbf{D}(I \times \bDelta_1)$ and $\mathbf{D}(J \times \bDelta_1)$ respectively.

I am going to be guided by the special case of a trivial derivator
\[ \mathbf{D} = \mathbf{D}_{\mathcal{M}}, \]
where $\mathcal{M}$ is a ``naked'' category of models $(\Sigma = \emptyset)$.

\msp{page 79}
For fixed $f$, we are going to construct an auxiliary category $K(f)$, in such a way that for every category $\mathcal{M}$, the systems $(G,F,u)$ consisting of an $\mathcal{M}$-presheaf $G$ on $J$, another $F$ on $I$, and an $f$-homomorphism $G \longrightarrow F$, are identified (functorially) with the $\mathcal{M}$-presheaves on
\printedpage{5}
% head: A. Grothendieck : Dérivateurs, Ch. IV, Édition des Universités de Montpellier II et Paris VII
% foot: Version 2 février 2002
$K(f)$. That said, we are going to try to interpret the skeletal arrows (1), but no matter here -- all the $\mathbf{D}$-skeletons are now fleshed out) in terms of the geometry of $K(f)$. % [sic] remplumés) -> remplumés, closing parenthesis printed without any opening parenthesis

\[
\left\{
\begin{aligned}
&\Ob K(f) = \Ob K(I) \amalg \Ob K(J) \qquad \text{\texttt{[rather} $\Ob I \amalg \Ob J$\texttt{ ]}}\\
&\text{$I$, $J$ being identified with full subcategories of $K(f)$,}\\
&\text{there are no arrows from a $j \in J$ to an $i \in I$.}\\
&\text{(Thus $I$ is open in $K(f)$, $J$ its closed complement.)}\\
&\text{For $i \in \Ob I$, $j \in \Ob J$, we set $\Hom_{K(f)}(i,j) = \Hom_J(f(i),j)$.}
\end{aligned}
\right.
\]

The composition of arrows in $K(f)$ is obvious. We are going to define a canonical retraction
\[ \bar{f} : K(f) \longrightarrow J \]
of $K(f)$ onto its closed subcategory $J$. Thus $\bar{f}$ is known on $\Ob J$, and on $\Ob I$ it will be $f$:
\[
\left\{
\begin{aligned}
\bar{f}(j) &= j && \text{if } j \in \Ob J\\
\bar{f}(i) &= f(i) && \text{if } i \in \Ob I.
\end{aligned}
\right.
\]
It remains to define, for $x, y \in \Ob K(f)$,
\msp{page 80}
\[ \Hom_{K(f)}(x,y) \longrightarrow \Hom_J(\bar{f}(x), \bar{f}(y)). \]
If $x, y \in \Ob J$, it is the identity, if $x, y \in I$, it is $f$. If $x \in I$, $y \in J$, we must define an arrow
\[ \Hom(f(x),y) \longrightarrow \Hom(f(x),y), \]
and we again take the identity. We have
\[
\begin{tikzcd}
J \arrow[r, hook, "j"] \arrow[dr, equal] & K(f) \arrow[d, "\bar{f}"] & I \arrow[l, hook', "i"'] \\
& J &
\end{tikzcd}
\]
\[
\left\{
\begin{aligned}
\bar{f} \circ j &= \id_J\\
j \circ \bar{f} &= \left\{
\begin{aligned}
&j && \text{on } J\\
&jf && \text{on } I,
\end{aligned}
\right.
\end{aligned}
\right.
\qquad \text{i.e.} \quad
\boxed{
\begin{aligned}
j\bar{f}j &= j && \text{on } J\\
j\bar{f}i &= jf && \text{on } I
\end{aligned}
}\ ,
\]
and we have \texttt{[}$K = K(f)$\texttt{]}
\[
\begin{aligned}
\alpha :\ \id_K &\longrightarrow j\bar{f} && \text{i.e. functorial homomorphism}\\
x &\xrightarrow{\alpha(x)} j(\bar{f}(x)), && x \in \Ob K,
\end{aligned}
\]
\printedpage{6}
% head: A. Grothendieck : Dérivateurs, Ch. IV, Édition des Universités de Montpellier II et Paris VII
% foot: Version 2 février 2002
defined as follows:
\[
\left\{
\begin{aligned}
&\text{If } x \in J,\ \alpha(x) = \id_x\ :\ x \xrightarrow{\sim} j(\bar{f}(x)) = j(x) = x\\
&\text{If } x \in I,\ \alpha(x) : \underbrace{x \longrightarrow f(x)}_{\text{in } K(f)},\ \text{or } \underbrace{f(x) \longrightarrow f(x)}_{\text{in } J},\ \alpha(x) = \id_{f(x)}.
\end{aligned}
\right.
\]
We must check functoriality for $x$ varying in $K$ \dots

Let $\Phi$ be an $\mathcal{M}$-presheaf on $K$. It defines
\[
\begin{aligned}
F &= i^*(\Phi) && \in \Ob \mathbf{D}_{\mathcal{M}}(I)\\
G &= j^*(\Phi) && \in \Ob \mathbf{D}_{\mathcal{M}}(J).
\end{aligned}
\]
On the other hand $f^*(G)$ is interpreted as $(j\bar{f})^*(\Phi)|I$, since
\[
\begin{aligned}
f^*(G)(i) &= G(\ \underbrace{f(i)}_{= \bar{f}(i) \,\in\, J}\ )\\
&= j^*(\Phi)(\bar{f}(i))\\
&= \Phi(j\bar{f}(i))\\
&= (j\bar{f})^*(\Phi)(i).
\end{aligned}
\]

\msp{page 81}
Thus
\[ f^*(G) \simeq (j\bar{f})^*(\Phi)|I. \]
Now the arrow
\[ \alpha : \id_K \longrightarrow j\bar{f} \]
gives
\[ \underbrace{\id_K(\Phi)}_{= \Phi} \xleftarrow{\Phi(\alpha)} (j\bar{f})^*(\Phi), \qquad \text{functorial in } \Phi, \]
whence, taking the restrictions to $I$,
\[ i^*(\Phi) \overset{\text{def}}{=} F \xleftarrow{\ u\ } f^*G. \]

In fact, without too much twisting and turning, we should be able to show that for every category $\mathcal{M}$, the category $\uHom(K^o, \mathcal{M})$ of $\mathcal{M}$-presheaves on $K$ is isomorphic to that of triples $(F,G,u)$, with $F \in \Ob \uHom(I^o, \mathcal{M})$, $G \in \uHom(J^o, \mathcal{M})$, and $u : f^*(G) \longrightarrow F$. Thus, $K$ appears as a solution of a universal problem
\[ \uHom(K^o, \mathcal{M}) \longrightarrow \{F,\ G,\ u \text{ as above}\}(\mathcal{M}). \]

But to specify the arrow $\longrightarrow$ here in a simple conceptual way, it is good to give ourselves $K$ with the additional structures \begin{tikzcd}[cramped, column sep=normal] I \arrow[r, hook, shift left, "j"] & K \arrow[l, shift left, "\bar{f}"] & I \arrow[l, hook', "i"'] \end{tikzcd}, $\bar{f}j = \id_J$, $\id_K \xrightarrow{\alpha} j\bar{f}$. And this system % [sic] I (first object, source of j and target of f-bar) -> J, since j : J -> K(f) and f-bar : K(f) -> J
\printedpage{7}
% head: A. Grothendieck : Dérivateurs, Ch. IV, Édition des Universités de Montpellier II et Paris VII
% foot: Version 2 février 2002
\msp{page 82}
$\{K,\ i,\ j,\ \bar{f},\ \alpha\}$ is (unless I am mistaken) defined up to canonical isomorphism by the condition of representing the functor $\mathcal{M} \longmapsto \{F,\ G,\ u\}^o(\mathcal{M})$ \dots

In particular, the system defined by $G \in \Ob \mathbf{D}_{\mathcal{M}}(J)$
\[ (f^*(G),\ G,\ \id : f^*(G) \longrightarrow f^*(G)) \]
is defined unless I am mistaken by $\bar{f}^*(G)$ on $K(f)$. We propose [?] to find a similar interpretation for
\[ (F,\ f_*(F),\ f^*f_*F \longrightarrow F) \]
$(F \in \Ob \mathbf{D}(I))$. The natural candidate here would be $\Phi = i_*(F)$. Let us grant this point (i.e. especially the verification of the fact that $j^*(i_*(F)) \simeq f_*(F)$). All this shows that if $\mathbf{D}$ is once again an arbitrary derivator, we can, by the preceding formulas, find a canonical object $\Phi$ of $\mathbf{D}(K)$, in terms either of a $\mathbf{D}$-coefficient $F$ on $I$ (taking $\Phi = i_*(F)$), or of a $\mathbf{D}$-coefficient $G$ on $J$ (taking $\Phi = \bar{f}^*(G)$). On the other hand, we can hope to define canonically
\[
\begin{aligned}
\mathbf{D}(K) &\longrightarrow \mathbf{D}(I \times \bDelta_1)\\
\mathbf{D}(K) &\longrightarrow \mathbf{D}(J \times \bDelta_1)
\end{aligned}
\]
whose composites with the ``skeleton'' functors
\[
\begin{aligned}
\mathbf{D}(I \times \bDelta_1) &\longrightarrow \uHom(\bDelta_1^o, \mathbf{D}(I))\\
\mathbf{D}(J \times \bDelta_1) &\longrightarrow \uHom(\bDelta_1^o, \mathbf{D}(J))
\end{aligned}
\]
\msp{page 83}
give back (for $\Phi \in \Ob \mathbf{D}(K)$) the skeletal arrows
\[
\left\{
\begin{aligned}
f^*(G) &\xrightarrow{\ u\ } F\\
G &\xrightarrow{\ v\ } f_*F
\end{aligned}
\right.
\]
associated with $(F,G,u)$. To be precise, $f^*(G) \longrightarrow F$ is defined in terms of $\mathbf{D}$, and $\Phi \in \Ob \mathbf{D}(K)$, as above, and $v$ is deduced from it by adjunction.

The natural idea here, would be to reduce the problem of fleshing out the skeletal adjunction arrows, to the case of open or closed immersions. Note that the open subcategories of $K$ correspond precisely to the functors $K \longrightarrow \bDelta_1$, by taking the inverse images of the generic open. There is therefore a unique % [sic] sqelettiques -> squelettiques, missing u (spelled squelettiques elsewhere in the text)
\[ \varphi :\ K(f) \longrightarrow \bDelta_1 \]
with $I = \varphi^{-1}(\{0\})$. Thus \texttt{[end of the chapter]}

\end{document}
